\section{Tipe Data: number, string, boolean, undefined, null}

JavaScript memiliki tipe data \textbf{primitif} dan \textbf{objek}. Tipe primitif menyimpan nilai tunggal dan bersifat immutable. Ada tujuh tipe primitif, lima di antaranya paling sering digunakan: \code{number}, \code{string}, \code{boolean}, \code{undefined}, dan \code{null}. Operator \code{typeof} mengembalikan string yang menunjukkan tipe data sebuah nilai \cite{jsInfo}.

\begin{table}[htbp]
\centering
\begin{tabular}{|P{2cm}|P{3cm}|P{3cm}|P{4cm}|}
\hline
\textbf{Tipe} & \textbf{Contoh} & \textbf{typeof} & \textbf{Keterangan} \\
\hline
number & \code{42}, \code{3.14}, \code{NaN} & \code{"number"} & Integer dan float; termasuk NaN, Infinity \\
\hline
string & \code{"Hello"}, \code{`Hi`} & \code{"string"} & Teks; kutip tunggal/ganda/backtick \\
\hline
boolean & \code{true}, \code{false} & \code{"boolean"} & Dua nilai logika \\
\hline
undefined & variabel tanpa nilai & \code{"undefined"} & Default variabel yang belum diisi \\
\hline
null & \code{null} & \code{"object"} & Nilai ``kosong'' eksplisit (bug typeof) \\
\hline
\end{tabular}
\caption{Lima tipe data primitif JavaScript yang paling umum}
\end{table}

JavaScript bersifat \textbf{dynamically typed}: tipe data variabel ditentukan oleh nilai yang disimpan, bukan oleh deklarasi. Variabel yang sama dapat menyimpan number, lalu diubah menjadi string. Meskipun fleksibel, ini dapat menyebabkan bug jika tidak berhati-hati. Template literal (backtick \code{`Hello \$\{name\}`}) memungkinkan interpolasi string yang lebih mudah dibaca dibanding konkatenasi \cite{mdnJsGuide}.

\begin{catatan}
\textbf{Nilai Truthy dan Falsy.} Dalam konteks boolean, beberapa nilai dianggap \textit{falsy} (dievaluasi sebagai \code{false}): \code{false}, \code{0}, \code{""} (string kosong), \code{null}, \code{undefined}, dan \code{NaN}. Semua nilai lain adalah \textit{truthy}. Ini penting untuk percabangan: \code{if (name)} akan gagal jika \code{name} adalah string kosong, null, atau undefined \cite{jsInfo}.
\end{catatan}

